\scanpage{38}
Soit $\rho\in\Gamma(S,\mathcal O_S)$, $U_\rho\subset E^1_S$ formé des $\lambda$ tels que
$1+\lambda\rho$ inversible i.e.
\[
U_\rho=\uline{\operatorname{Spec}}\ \mathcal O_S[T]_{(1+\rho T)}
\]
prenons la multiplication
\[
\unsure{(}\lambda *_\rho\lambda'=\lambda+\lambda'+\lambda\lambda'\rho
\]
de sorte que $U_\rho$ est un schéma en groupes
\inserted{(\unsure{Assoc.} comm. \unsure{él. unit.} $\unsure{1}$)}, et
\diagram{\begin{tikzcd}[ampersand replacement=\&,column sep=large,row sep=small]
\lambda \arrow[r,maps to,no head] \& 1+\lambda\rho\\
U_\rho \arrow[r,no head] \& \mathbb G_m
\end{tikzcd}}
est un hom. de groupes, qui est un iso. sss $\rho$ inversible.

Considérons un hyperpolyèdre régulier \unsure{ponctué|pointé} épinglé sur $S$, de dim. $n$,
\inserted{(\unsure{conn.})}
correspondant à des invariants $\mu_i$ $(0\le i\le n)$ et $\beta_i$
$(0\le i\le n-1)$ [avec $\mu_i+1$ inv. $0\le i\le n$).
\unsure{Soit}
\deleted{$\rho=\unsure{\mathbb L}_{n+1}(\mu_0,\lambda_0,\ldots,\unsure{\lambda_{n-1}},\mu_n)$}
\[
\unsure{=}\qquad
\rho=\unsure{\mathbb L}_{n+1}(\mu_0,\lambda_0,\ldots,\lambda_{n-1},\mu_n)
=\langle\Pi,\Pi'\rangle
\]
\mspage{37}

